\documentclass[11pt]{article}
\usepackage[margin=1in]{geometry}
\usepackage{amsmath,amssymb,amsthm,mathtools}
\usepackage{graphicx}
\usepackage{microtype}
\usepackage[hidelinks]{hyperref}

\newcommand{\R}{\mathbb R}
\newcommand{\cE}{\mathcal E}
\newtheorem{theorem}{Theorem}
\newtheorem{corollary}{Corollary}

\title{A purely metric parallel-sum formula\\
for infimal convolution on geodesic products}
\author{Candidate full answer for arXiv:2503.12939}
\date{19 July 2026}

\begin{document}
\maketitle

\begin{center}
\fbox{\parbox{0.91\textwidth}{
\textbf{Status: candidate full answer, likely valid.}
Weighted product metrics on arbitrary finite products of geodesic spaces have
an explicit metric infimal convolution, which is again a finite nondegenerate
geodesic distance.  This gives a broad affirmative answer to the source's
existential question, subject to human review.}}
\end{center}

\section{Source question}

The source is Nicolò De Ponti, Giacomo Enrico Sodini, and Luca Tamanini,
\emph{The infimal convolution structure of the Hellinger-Kantorovich
distance}, arXiv:2503.12939v1 \cite{DPST}.  After showing that metric infimal
convolution can be degenerate or fail the triangle inequality, and after
discussing the more structured normed-space case, the authors ask in the
Introduction (page 5):

\begin{quote}
``This raises the compelling question of whether there exist other purely
metric examples, similar to the Hellinger--Wasserstein case, where this
construction yields meaningful distances.''
\end{quote}

\begin{figure}[ht]
\centering
\includegraphics[width=\textwidth]{figures/open_problem_crop.png}
\caption{The source question in its full page-width context, page 5.}
\label{fig:question}
\end{figure}

For two costs \(c_1,c_2\) on a set \(Z\), the source defines the square of
their metric infimal convolution at \(z^0,z^1\) by
\begin{equation}\label{eq:def}
(c_1\nabla c_2)^2(z^0,z^1)
=\liminf_{N\to\infty}\inf_P
N\sum_{i=1}^N\left(c_1^2(x_{i-1},y_i)+c_2^2(y_i,x_i)\right),
\end{equation}
where \(P=(x_0,\ldots,x_N;y_1,\ldots,y_N)\), \(x_0=z^0\), and
\(x_N=z^1\).

\section{Proof intuition}

On a product of geodesic spaces, give each factor one weight in the first
metric and another weight in the second.  During an alternating step, each
coordinate can place its intermediate point at its own optimal fraction along
a geodesic.  Minimizing the two weighted squared lengths produces the familiar
parallel-sum coefficient
\[
\left(a_1^{-2}+a_2^{-2}\right)^{-1/2}.
\]
The same elementary quadratic inequality gives a lower bound for every
possible path.  Thus the effective product distance is not merely the limit:
its square is the exact minimum energy at every discretization level \(N\).

\section{The product theorem}

\begin{theorem}[Purely metric parallel sum]\label{thm:main}
Let \((X_k,d_k)\), \(k=1,\ldots,m\), be geodesic metric spaces and let
\(Z=\prod_{k=1}^m X_k\).  For positive numbers \(a_{jk}\), \(j=1,2\),
define metrics
\begin{equation}\label{eq:inputs}
\rho_j(z,z')^2=\sum_{k=1}^m a_{jk}^2d_k(z_k,z'_k)^2,
\qquad j=1,2.
\end{equation}
Set
\begin{equation}\label{eq:weights}
a_{*k}:=\left(a_{1k}^{-2}+a_{2k}^{-2}\right)^{-1/2}
=\frac{a_{1k}a_{2k}}{\sqrt{a_{1k}^2+a_{2k}^2}}
\end{equation}
and
\begin{equation}\label{eq:output}
\rho_*(z,z')^2:=\sum_{k=1}^m a_{*k}^2d_k(z_k,z'_k)^2.
\end{equation}
Then
\[
\rho_1\nabla\rho_2=\rho_*.
\]
More precisely, for every \(N\ge1\) and every pair of endpoints, the infimum
of the \(N\)-path energy in \eqref{eq:def} is exactly the square of their
\(\rho_*\)-distance.
\end{theorem}

\begin{proof}
Each function in \eqref{eq:inputs} and \eqref{eq:output} is a metric.  This
follows from the coordinate triangle inequalities and Minkowski's inequality.
They are geodesic: choose constant-speed geodesics in every factor and run
them with a common time parameter.

Fix endpoints \(z^0,z^1\) and an arbitrary \(N\)-path
\(P=(x_0,\ldots,x_N;y_1,\ldots,y_N)\) between them.  For each \(i,k\), put
\[
A_{ik}=d_k(x_{i-1,k},y_{i,k}),\quad
B_{ik}=d_k(y_{i,k},x_{i,k}),\quad
C_{ik}=d_k(x_{i-1,k},x_{i,k}).
\]
The triangle inequality gives \(C_{ik}\le A_{ik}+B_{ik}\).  For positive
\(a,b\) and nonnegative \(A,B\),
\begin{equation}\label{eq:scalar}
a^2A^2+b^2B^2
-\frac{a^2b^2}{a^2+b^2}(A+B)^2
=\frac{(a^2A-b^2B)^2}{a^2+b^2}\ge0.
\end{equation}
Applying \eqref{eq:scalar} coordinatewise yields
\begin{align*}
\rho_1^2(x_{i-1},y_i)+\rho_2^2(y_i,x_i)
&=\sum_k\left(a_{1k}^2A_{ik}^2+a_{2k}^2B_{ik}^2\right)\\
&\ge\sum_k a_{*k}^2(A_{ik}+B_{ik})^2\\
&\ge\sum_k a_{*k}^2C_{ik}^2
=\rho_*^2(x_{i-1},x_i).
\end{align*}
Consequently, Cauchy--Schwarz and the triangle inequality for \(\rho_*\) give
\begin{align}\label{eq:lower}
\cE_N(P)
&\ge N\sum_{i=1}^N\rho_*^2(x_{i-1},x_i)\\
&\ge\left(\sum_{i=1}^N\rho_*(x_{i-1},x_i)\right)^2
\ge\rho_*^2(z^0,z^1).\nonumber
\end{align}

It remains to attain the lower bound.  For each coordinate choose a
constant-speed geodesic \(\gamma_k:[0,1]\to X_k\) from \(z_k^0\) to
\(z_k^1\), and write \(D_k=d_k(z_k^0,z_k^1)\).  Define
\begin{equation}\label{eq:path}
x_{i,k}=\gamma_k(i/N),\qquad
t_k=\frac{a_{2k}^2}{a_{1k}^2+a_{2k}^2},\qquad
y_{i,k}=\gamma_k\left(\frac{i-1+t_k}{N}\right).
\end{equation}
Then
\[
A_{ik}=\frac{t_kD_k}{N},\qquad
B_{ik}=\frac{(1-t_k)D_k}{N}.
\]
The choice of \(t_k\) is the equality case in \eqref{eq:scalar}, so every
alternating step in \eqref{eq:path} has cost
\[
\sum_k\left(a_{1k}^2A_{ik}^2+a_{2k}^2B_{ik}^2\right)
=\frac1{N^2}\sum_k a_{*k}^2D_k^2
=\frac{\rho_*^2(z^0,z^1)}{N^2}.
\]
Multiplying the sum of the \(N\) step costs by \(N\) gives
\(\cE_N(P)=\rho_*^2(z^0,z^1)\).  Together with \eqref{eq:lower}, this proves
the exact minimum formula for every \(N\), and hence the asserted infimal
convolution identity.
\end{proof}

\section{A branching, non-vector example}

Let \(T\) be the metric tripod obtained by gluing three copies of \([0,1]\)
at their zero endpoints, with path distance \(d_T\).  On
\(Z=T\times[0,1]\), define
\begin{align*}
\rho_1((x,s),(y,t))^2&=d_T(x,y)^2+4|s-t|^2,\\
\rho_2((x,s),(y,t))^2&=4d_T(x,y)^2+|s-t|^2.
\end{align*}
In each coordinate the effective weight is
\[
\left(1^{-2}+2^{-2}\right)^{-1/2}=\frac2{\sqrt5}.
\]
Theorem~\ref{thm:main} therefore gives
\begin{equation}\label{eq:tripod}
(\rho_1\nabla\rho_2)((x,s),(y,t))
=\frac2{\sqrt5}\sqrt{d_T(x,y)^2+|s-t|^2}.
\end{equation}
This is a finite nondegenerate geodesic distance on a space with genuine
branching.  The inputs are not proportional: on tripod fibers \(\rho_2=2\rho_1\),
whereas on interval fibers \(\rho_1=2\rho_2\).  Thus both inputs contribute
nontrivially to the output.

\section{Verification, scope, and novelty}

The included exact checker expands \eqref{eq:scalar}, verifies the optimizing
fraction and effective coefficient symbolically, checks the tripod constants,
and tests the lower bound on deterministic rational product paths.  Run
\begin{verbatim}
conda run --no-capture-output -n sandbox python \
  code/check_parallel_sum.py
\end{verbatim}
from the packet directory.  Computation is only an audit; the proof above is
exact.

The result answers the source's existential question affirmatively with a
large nonlinear class.  It does not classify every well-behaved pair, prove
that arbitrary geodesic metrics have a distance-valued infimal convolution,
or address the separate question about \(p\)-Hellinger--Kantorovich distances.

Before promotion, the run indexes were searched using arXiv:2503.12939, the
source title, the exact question terms, metric infimal convolution, geodesic
products, weighted product metrics, and parallel sums.  A bounded arXiv/web
search used the exact source question and close product/geodesic/parallel-sum
phrases.  It found the source paper and unrelated functional
infimal-convolution literature, but no separate statement of
Theorem~\ref{thm:main}.  The current arXiv record is v1, submitted 17 March
2025, and still contains the question.  This records the search bounds and is
not an exhaustive novelty certification.

Human review is recommended.  The main scope question is whether this broad
existential answer matches what the source authors intend by ``purely metric'';
mathematically, the highest-value checks are the scalar equality case and the
coordinatewise attaining path \eqref{eq:path}.

\begin{thebibliography}{9}
\bibitem{DPST}
N. De Ponti, G. E. Sodini, and L. Tamanini,
\emph{The infimal convolution structure of the Hellinger-Kantorovich
distance}, arXiv:2503.12939v1, 2025.
\url{https://arxiv.org/abs/2503.12939}.
\end{thebibliography}

\end{document}
